%% ===================================================================
%%  Plantilla lista para usar -- formato Wordglow
%%  Compilar:  pdflatex plantilla.tex   (dos veces)
%%  Referencia de todos los elementos:  ejemplo.tex
%% ===================================================================
\documentclass[letterpaper]{wordglow}   % o [a4paper]

%% ---- Datos del documento -------------------------------------------
%% El titulo se escribe TAL CUAL sale impreso.  \\ corta linea.
%% \wgac{...} lleva esa palabra al color de acento.
\wgTitulo{DOCUMENTACI\'ON\\\wgac{T\'ECNICA} DEL\\PROYECTO}
\wgEyebrow{Desarrollo de software}
\wgKicker{Informe final}{Dise\~no y desarrollo}
\wgAutor{Nombre Apellido}
\wgVersion{1.0}
\wgFecha{\today}

%% ---- Marca (encabezado, pie y pagina de cierre) -------------------
\wgMarca{Wordglow}
\wgSubtitulo{Documentaci\'on t\'ecnica}
\wgWeb{www.wordglow.com}
\wgCorreo{contacto@wordglow.com}
\wgTel{+591 000 0000}
\wgDireccion{Av. Ejemplo 123, Cochabamba, Bolivia}

\begin{document}

\maketitle

\pagenumbering{roman}
\wgindice
\clearpage
\pagenumbering{arabic}

% ===================================================================
\section{T\'itulo de la primera secci\'on}

Escribe aqu\'i el contenido. El cuerpo, el interlineado y los m\'argenes
ya est\'an fijados por la clase.

\subsection{Subsecci\'on}

M\'as texto.

% -- Elementos disponibles (ver ejemplo.tex para todos):
% \wgchip{Secci\'on 1}                      chip enmarcado sobre un titulo
% \wgnota{...}                              callout con regla de acento
% \begin{quote}...\wgfuente{Autor}\end{quote}   cita
% \begin{lstlisting}...\end{lstlisting}     bloque de codigo
% \wgcodigotitulo{...} + lstlisting         codigo con rotulo "Codigo N"
% \wgecuacion{formula}{n}                   ecuacion numerada entre reglas
% figure + \caption + \wgplaceholder + \wgnotafig{...}
% table  + \caption + tabular(L{}/R{}) + \wgth/\wgtoprule/\wgrowrule

% ===================================================================
\clearpage
\begin{referencias}
\refitem Apellido, A.~A. (2026). \textit{T\'itulo de la obra}. Editorial.
\end{referencias}

\wgcierre

\end{document}
